% Демонстрация 6: «Корни многочлена и теорема Виета» (Лекция 13)
\section*{Демонстрация 6. «Корни многочлена и теорема Виета»\\ \normalsize\textmd{Файл \texttt{demos/06-polynomial-vieta.html} \quad(к Лекции~13)}}

Демонстрация представляет собой интерактивную модель: на комплексной
плоскости изображены корни $x_1, \dots, x_n$ многочлена, которые
слушатель перетаскивает мышью. Комплексные корни перемещаются
\emph{парами сопряжённых} (ведомый корень автоматически отражается
относительно действительной оси), действительные корни скользят вдоль
оси~$\operatorname{Re}$. По заданным корням программа в реальном
времени пересобирает многочлен
\[
  P(x) = a_n \prod_{k=1}^{n} (x - x_k),
\]
выводит его коэффициенты в развёрнутом виде, проверяет соотношения
теоремы Виета и строит график $P(x)$ на действительной оси. Ведущий
коэффициент $a_n$ регулируется отдельным ползунком.

\subsection*{Какие факты лекции иллюстрирует}

Демонстрация даёт наглядную, «осязаемую» опору сразу для нескольких
ключевых результатов Лекции~13.

\begin{itemize}
  \item \textbf{Разложение на линейные множители}
    (теорема о разложении многочлена на линейные множители,
    $P(x)=a_n\prod_{k=1}^{n}(x-x_k)$). Слушатель буквально
    \emph{задаёт} корни $x_1, \dots, x_n$, а программа собирает
    \[
      P(x) = a_n (x - x_1)(x - x_2)\cdots(x - x_n).
    \]
    Перетаскивая корень, студент видит, как одновременно меняется
    левая часть (множители) и правая (коэффициенты) — то есть как
    корни \emph{порождают} многочлен, а не наоборот.

  \item \textbf{Связь корней и коэффициентов — теорема Виета}
    (теорема о Виете в лекции). Панель «Теорема Виета» в каждый момент
    показывает проверку соотношений
    \[
      \sum_{k} x_k = -\frac{a_{n-1}}{a_n}, \qquad
      \prod_{k} x_k = (-1)^n \frac{a_0}{a_n},
    \]
    отмечая совпадение значений галочкой. Это превращает абстрактные
    формулы Виета в наблюдаемый, проверяемый на глазах факт.

  \item \textbf{Кратность корня} (определение кратности и
    критерий кратности). Сливая два действительных корня в
    одну точку, студент получает корень кратности~$2$ и видит на
    графике, как линия перестаёт \emph{пересекать} ось и начинает её
    \emph{касаться}. Это прямая геометрическая интерпретация
    множителя $(x - \alpha)^k$.

  \item \textbf{Комплексно-сопряжённые корни} многочленов с
    действительными коэффициентами (теорема о сопряжённых корнях).
    Поскольку в демонстрации комплексные корни всегда живут парами
    $x$ и $\bar{x}$, коэффициенты собранного многочлена остаются
    действительными — что и видно по выводимой формуле. Это
    «физически» воспроизводит вывод лекции: пара
    $(x - \alpha)(x - \bar\alpha) = x^2 + px + q$ с действительными
    $p, q$.

  \item \textbf{Теорема Безу и остаток.} Хотя демонстрация не делит
    многочлены явно, она удобна как мост к теореме Безу: каждый
    отмеченный на графике действительный корень — это точка, где
    $P(x) = 0$, то есть остаток от деления на $(x - \alpha)$ обращается
    в нуль (теорема Безу и её следствие).
\end{itemize}

\begin{remarkIT}{Зачем это нужно на практике}{demo06_apps}
  Связь «корни\,$\leftrightarrow$\,коэффициенты», которую исследует
  демонстрация, лежит в основе целого ряда инженерных задач.
  \begin{itemize}
    \item \textbf{Устойчивость систем.} Поведение линейной системы
      управления определяется корнями её характеристического
      многочлена: система устойчива, когда все корни лежат в левой
      полуплоскости $\operatorname{Re} x < 0$. Перетаскивая корень
      через мнимую ось, можно показать «границу устойчивости».
    \item \textbf{Цифровые фильтры (DSP).} Нули и полюса передаточной
      функции — это корни числителя и знаменателя. Их расположение на
      плоскости задаёт амплитудно-частотную характеристику фильтра, а
      сопряжённые пары объединяются в действительные звенья второго
      порядка (см.\ замечание о biquad-секциях в лекции).
    \item \textbf{Безопасность вычислений.} Кратные и близкие корни
      численно неустойчивы: малое возмущение коэффициентов сильно
      смещает корни. Демонстрация наглядно показывает, почему рядом с
      кратным корнем график «прижимается» к оси.
    \item \textbf{Символьные вычисления (CAS).} Системы вроде
      \texttt{SymPy}, \texttt{Mathematica}, \texttt{Maxima} переходят
      между формами $a_n\prod(x - x_k)$ и списком коэффициентов
      именно по формулам Виета — той самой операцией, которую
      выполняет демонстрация.
  \end{itemize}
\end{remarkIT}

\subsection*{Примерный сценарий использования на занятии}

Демонстрацию удобно включать сразу после доказательства разложения на
линейные множители $P(x)=a_n\prod_{k=1}^{n}(x-x_k)$ и формулировки теоремы
Виета — когда формулы уже выписаны, но ещё не «прожиты». Ниже —
примерный сценарий на 7--10 минут.

\begin{enumerate}
  \item \textbf{Старт от готовой картинки.} Откройте демонстрацию с
    начальной конфигурацией (два действительных корня и одна пара
    сопряжённых). Покажите три согласованных представления одного
    объекта: корни на плоскости, формулу
    $P(x) = a_n\prod(x - x_k)$ и развёрнутый многочлен с
    коэффициентами. Подчеркните, что это \emph{один и тот же}
    многочлен в разных формах.

  \item \textbf{Двигаем действительный корень.} Перетащите красный
    корень вдоль оси и попросите следить за графиком: точка
    пересечения с осью~$x$ едет вместе с корнем. Связь «корень — нуль
    функции» (теорема Безу) становится очевидной.

  \item \textbf{Выносим корень в комплексную область.} Возьмите
    действительный корень (или ведущий корень пары) и утащите его
    вверх от оси. \emph{Момент для вопроса:}
    \begin{quote}
      «Смотрите — как только корень покинул действительную ось, тут же
      появился его двойник снизу. Почему программа не позволяет
      задать \emph{одиночный} комплексный корень, оставив
      коэффициенты действительными?»
    \end{quote}
    Подведите к ответу: по теореме о сопряжённых корнях у многочлена с
    действительными коэффициентами комплексные корни обязаны идти
    парами $x, \bar x$. Обратите внимание класса, что при этом
    выводимые коэффициенты остаются \emph{действительными числами} —
    мнимые части сократились.

  \item \textbf{Проверяем Виета вживую.} Не закрывая панель «Теорема
    Виета», медленно подвигайте любой корень и обратите внимание, что
    сумма $\sum x_k$ и произведение $\prod x_k$ всё время совпадают с
    $-a_{n-1}/a_n$ и $(-1)^n a_0/a_n$ (галочки горят). Здесь уместно
    устно решить обратную задачу: «Если мы хотим сумму корней,
    равную~$0$, как нужно расположить корни?»

  \item \textbf{Сливаем корни — момент удивления.} Подведите два
    действительных корня вплотную друг к другу. \emph{Кульминация
    демонстрации:} в точке слияния график \emph{перестаёт пересекать}
    ось и лишь \emph{касается} её.
    \begin{quote}
      «Что произошло с графиком? Почему он больше не пересекает ось, а
      только дотрагивается до неё?»
    \end{quote}
    Свяжите наблюдение с множителем $(x - \alpha)^2$ и определением
    кратности: касание оси — геометрический признак кратного корня.

  \item \textbf{Меняем ведущий коэффициент.} Подвигайте ползунок
    $a_n$. Корни на месте — а график растягивается/отражается. Вывод:
    корни задают \emph{где} нули, а $a_n$ — только общий масштаб
    (и знак) многочлена. Это хорошо согласуется с тем, что в формулах
    Виета всюду стоит отношение $a_k/a_n$.

  \item \textbf{Степень и число корней.} Добавьте ещё одну
    сопряжённую пару и обратите внимание на счётчик степени: степень
    всегда равна числу корней с учётом кратности — иллюстрация
    основной теоремы алгебры.

  \item \textbf{Связка с задачами лекции.} Завершите переходом к
    разобранным в лекции примерам: «Мы только что \emph{собирали}
    многочлен из корней. На практике чаще приходится решать обратную
    задачу — по коэффициентам \emph{найти} корни. Этим и займёмся,
    используя схему Горнера и сопряжённость комплексных корней».
\end{enumerate}

\begin{remarkIT}{Методические акценты}{demo06_tips}
  \begin{itemize}
    \item Главные два «ага-момента» — автоматическое появление
      сопряжённого корня (действительность коэффициентов) и переход
      «пересечение\,$\to$\,касание» при слиянии корней (кратность).
      Не проговаривайте их заранее: дайте студентам сначала
      \emph{увидеть}, а затем сформулировать самим.
    \item Демонстрация показывает только действительные корни на
      графике $P(x)$; комплексные корни на графике «не видны». Это сам
      по себе полезный вопрос классу: «Почему у многочлена 4-й степени
      график может вообще не пересекать ось?»
    \item Демонстрация иллюстративна, а не доказательна: она
      \emph{подтверждает} теоремы на примерах, но не заменяет их
      доказательств, проведённых в лекции.
  \end{itemize}
\end{remarkIT}
